\section{Sonstiges}
\subsection{Beispiel gitternetz}\label{chap:bsp_gitternetz}

Die Funktion bildet die $z$-Ebene auf die $w$-Ebene ab ($w = w_1 + \jimg w_2$).
Mit der Zerlegung $z = z_1 + \jimg z_2$ ist die Transformation definiert als:
\[ w = (z_1 + \jimg z_2)^2 = (z_1^2 - z_2^2) + \jimg(2 z_1 z_2) \]
Daraus folgen die kartesischen Koordinaten in der $w$-Ebene:
\[
\begin{cases}
    w_1 = z_1^2 - z_2^2 \\
    w_2 = 2 z_1 z_2
\end{cases}
\]

Nun analysieren wir die Transformation, indem wir die Formeln der kartesischen Geraden in $f$ einsetzen:

\begin{minipage}[t]{0.48\linewidth}
    \textbf{Vertikale ($z = c_1 + \jimg s$):}
    Festhalten des Realteils $z_1 = c_1$:
    \[ \begin{cases} w_1 = c_1^2 - s^2 \\ w_2 = 2c_1s \rightarrow s = \frac{w_2}{2c_1} \end{cases} \]
    Einsetzen von $s$ in $w_1$:
    \[ \colorbox{blue!10}{$w_1 = c_1^2 - \frac{w_2^2}{4c_1^2}$} \]
    Stellt \textbf{Parabeln} mit der $w_1$-Achse dar, die nach links geöffnet sind.
\end{minipage}
\hfill
\begin{minipage}[t]{0.48\linewidth}
    \textbf{Horizontale ($z = r + \jimg c_2$):}
    Festhalten des Imaginärteils $z_2 = c_2$:
    \[ \begin{cases} w_1 = r^2 - c_2^2 \\ w_2 = 2rc_2 \rightarrow r = \frac{w_2}{2c_2} \end{cases} \]
    Einsetzen von $r$ in $w_1$:
    \[ \colorbox{green!10}{$w_1 = \frac{w_2^2}{4c_2^2} - c_2^2$} \]
    Stellt \textbf{Parabeln} mit der $w_1$-Achse dar, die nach rechts geöffnet sind.
\end{minipage}

\vspace{1.5em}

\subsubsection{Umgekehrtes Vorgehen}
\paragraph{Beispiel: Quadratisches Gitter}
\[f(z) = z^2 \quad \rightarrow \quad f(z) = (z_1 + z_2 \jimg)^2 = z_1^2 + \jimg 2 z_1
z_2 - z_2^2\]
\[w_1 = \RE(w) = z_1^2 - z_2^2 \qquad w_2 = \IM(w) = 2z_1z_2\]
Wir wissen, dass in einem quadr. Gitter $w_1 / w_2$ konstant sind ($x,y$), also:\vspace{1em}

\begin{minipage}{0.45\columnwidth}
    \[z_1^2 - z_2^2 = c_1\]
    \begin{center}
        Horizontale Hyperbel
    \end{center}
\end{minipage}\hfill
\begin{minipage}{0.45\columnwidth}
    \[c_2 = 2z_+z_2 \rightarrow z_2 = \frac{c_2}{2 z_1}\]
    \begin{center}
        Hyperbel $x^{-1}$
    
\end{center}
\end{minipage}

\subsection{Indexwechsel}\vspace{-1em}
\[ n = (\text{Schritt} \cdot m) + \text{Offset} \]

wobei $m = 0, 1, 2, \dots$ der neue Teilmatrix-Index ist (\textbf{m startet bei 0!}).
\begin{ctabular}{@{}lccccl@{}}
\toprule
\textbf{Ziel (n)} & \textbf{Schritt} & \textbf{Offset} & \textbf{Formel} & \textbf{Beispiel-Sequenza} \\ \midrule
Ungerade             & 2              & 1               & $n = 2m + 1$     & $1, 3, 5, 7, \dots$       \\
Gerade                & 2              & 2               & $n = 2m + 2$     & $2, 4, 6, 8, \dots$       \\
Alle 4 (ab 1)       & 4              & 1               & $n = 4m + 1$     & $1, 5, 9, 13, \dots$      \\
Alle 4 (ab 3)       & 4              & 3               & $n = 4m + 3$     & $3, 7, 11, 15, \dots$     \\ \bottomrule
\end{ctabular}

\newpage